\chapter{Reproducibility}
\label{app:reproducibility}

Every numerical result can be reproduced from the original CSV downloads
in approximately 45--60 minutes on a 2024 MacBook Pro M-series:

\begin{verbatim}
git clone <repository URL>
cd bsc-thesis-ids
python -m venv .venv && source .venv/bin/activate
pip install -r requirements.txt
# Place CICIDS2017 CSVs in data/raw/cicids2017/
# Place UNSW-NB15 CSVs in data/unsw-nb15/raw/
caffeinate -i make full     # macOS; on Linux use systemd-inhibit
\end{verbatim}

The \texttt{make full} target chains every step of the pipeline: data
preparation, sanity checks, training under both split policies, binary
training, evaluation, LODO cross-validation, UNSW training and
evaluation, MITRE mapping, traffic analysis, figure generation,
inference benchmarking, bootstrap confidence intervals, operating
points, anomaly detection, and adversarial evaluation. Individual
targets can be run in isolation.

The random seed is set to 42 throughout. Tree ensembles trained with
\texttt{n\_jobs=-1} introduce approximately 1\% F1 drift between runs
because of non-deterministic thread scheduling at leaf-split time.
Numbers in this thesis come from a single canonical run validated
against a second run to confirm the qualitative conclusions hold under
the drift.

\paragraph{Near-duplicate filtering sensitivity.}
A sensitivity study confirms that macro F1 on the binary day task
changes by less than 0.01 across all four models when near-duplicates
are filtered. This result should be read alongside the leakage-check
report, which flags many near-duplicate groups that cross the
train/test boundary in both stratified and day splits. Exact duplicates
are removed in cleaning; near-duplicates that survive are mostly
within-class neighbours (similar attack flows from the same campaign)
rather than benign-versus-attack confusables, so they do not materially
shift macro F1 on the binary task. A stricter near-duplicate pruning
would likely change per-class precision/recall and is the kind of
signal that Pendlebury et al.~\cite{Pendlebury2019} argue should be
reported explicitly rather than averaged away.

\begin{figure}[H]
    \centering
    \includegraphics[width=0.8\textwidth]{figure_V3_binary_recall_vs_fpr.png}
    \caption{Recall-versus-FPR sweep for all four binary models on the
    day test set. The full operating-point curve from which the headline
    operating points in Table~\ref{tab:operating-points} are drawn.}
    \label{fig:v3}
\end{figure}
